% ============================================================
%  Book chapter 20: Files
% ============================================================
\chapter{Files}

\textbf{What this chapter gives you:} how to open, read, write and close files safely, how to jump
around inside a file, and how to find, copy, move and name files and folders. You will use this
for JSONL eval sets, logs, model outputs, checkpoints and config files. You also meet
\code{StringIO} and \code{BytesIO}, which make memory behave like a file.

\tikzset{c20n/.style={draw=#1,fill=#1!10,rounded corners=3pt,minimum height=8mm,inner sep=4pt,
    font=\ttfamily\small},
  c20c/.style={draw=cgray,minimum width=7mm,minimum height=7mm,inner sep=0pt,font=\ttfamily\small}}

% ------------------------------------------------------------
\section{Files and directories}

A \emph{file} is a sequence of bytes stored under a name. A \emph{directory} (folder) holds files
and other directories, so the whole thing forms a tree. Text and bytes are the same on disk. The
only difference is whether Python decodes the bytes into \code{str} for you.

\begin{center}
\begin{tikzpicture}[>=Stealth,level distance=9mm,
    every node/.style={font=\ttfamily\small},
    d/.style={draw=cblue,fill=cblue!10,rounded corners=2pt,inner sep=3pt},
    f/.style={draw=cgray,fill=white,rounded corners=2pt,inner sep=3pt}]
  \node[d] (root) at (0,0) {project/};
  \node[d] (data) at (3.2,0.9) {data/};
  \node[d] (runs) at (3.2,-0.9) {runs/};
  \node[f] (a) at (6.6,1.5) {train.jsonl};
  \node[f] (b) at (6.6,0.7) {eval.jsonl};
  \node[f] (c) at (6.6,-0.5) {run-001.log};
  \node[f] (e) at (6.6,-1.3) {run-002.log};
  \draw[cgray] (root.east) -- (data.west);
  \draw[cgray] (root.east) -- (runs.west);
  \draw[cgray] (data.east) -- (a.west);
  \draw[cgray] (data.east) -- (b.west);
  \draw[cgray] (runs.east) -- (c.west);
  \draw[cgray] (runs.east) -- (e.west);
  \node[font=\sffamily\scriptsize,text=cblue,anchor=west] at (8.3,0.9) {directories};
  \node[font=\sffamily\scriptsize,text=cgray,anchor=west] at (8.3,-0.5) {files (leaves)};
\end{tikzpicture}
\end{center}
% MD-BEGIN
% ```text
% project/
% ├── data/
% │   ├── train.jsonl
% │   └── eval.jsonl
% └── runs/
%     ├── run-001.log
%     └── run-002.log
% ```
% MD-END

% ------------------------------------------------------------
\section{File Input and Output}

\subsection{Create or open with open()}

\code{open(name, mode)} returns a \emph{file object}. It is your handle on the file: you call
methods on it to read and write. When you are done you must close it, so that buffered writes
reach the disk.

\begin{center}
\begin{tabularx}{\linewidth}{@{}l X l l l@{}}
\toprule
\textbf{Mode} & \textbf{Meaning} & \textbf{File must exist?} & \textbf{Erases old content?} & \textbf{Pointer starts} \\
\midrule
\code{"r"} & read (default) & yes & no & start \\
\code{"w"} & write & no, created & \textbf{yes} & start \\
\code{"x"} & write, but only if new & must \emph{not} exist & -- & start \\
\code{"a"} & append & no, created & no & end \\
\code{"r+"} & read and write & yes & no & start \\
\bottomrule
\end{tabularx}
\end{center}

A second letter picks the kind: \code{t} for text (the default) or \code{b} for bytes. So
\code{"wb"} is write-bytes and \code{"rt"} is the same as \code{"r"}.

\begin{center}
\begin{tikzpicture}[>=Stealth]
  \node[c20n=cgreen,align=center] (m) at (0,0) {Python\\[-1pt]{\scriptsize\sffamily str in memory}};
  \node[c20n=cgray,align=center] (d) at (7.5,0) {file on disk\\[-1pt]{\scriptsize\sffamily bytes}};
  \draw[->,thick] ([yshift=2mm]m.east) -- node[above,font=\ttfamily\scriptsize]{write() = encode} ([yshift=2mm]d.west);
  \draw[->,thick] ([yshift=-2mm]d.west) -- node[below,font=\ttfamily\scriptsize]{read() = decode} ([yshift=-2mm]m.east);
  \node[font=\sffamily\scriptsize,text=cgray] at (3.75,-1.0) {text mode does the encoding for you; binary mode skips it};
\end{tikzpicture}
\end{center}
% MD-BEGIN
% ```mermaid
% flowchart LR
%     M["Python str (memory)"] -- "write() = encode" --> D["file on disk (bytes)"]
%     D -- "read() = decode" --> M
% ```
% *Text mode does the encoding for you; binary mode skips it.*
% MD-END

\begin{pycode}
f = open("notes.txt", "w", encoding="utf-8", newline="\n")
print(type(f))
n = f.write("line one\nline two\nline three\n")
print(n)
f.close()
print(f.closed)
\end{pycode}
\begin{rawcode}
<class '_io.TextIOWrapper'>
29
True
\end{rawcode}

\begin{typecheck}
\code{open(..., "w")} returns a \code{TextIOWrapper}, and \code{open(..., "wb")} returns a
\code{BufferedWriter}. They are not strings. \code{write()} returns how many characters (text
mode) or bytes (binary mode) it wrote, and it adds nothing: no space, no newline. A closed
file object still exists, but using it raises \code{ValueError}.
\end{typecheck}

\begin{gamerule}
Always pass \code{encoding="utf-8"} in text mode. If you leave it out, Python uses the system
default, which on many Windows machines is a legacy code page, so the same script can write
different bytes on different computers.
\end{gamerule}

\begin{trap}
Mode \code{"w"} destroys the old content the moment you open the file, before you write anything.
A typo in the filename can wipe an important file. Use \code{"x"} when the file must be new, and
\code{"a"} to add to the end (for example, a log or a JSONL results file).
\end{trap}

\subsection{with: close files automatically}

Writing \code{f.close()} by hand is easy to forget, and it is skipped if an error happens in the
middle. A \code{with} block closes the file for you when the block ends, normally or by an
exception.

\begin{center}
\begin{tikzpicture}[>=Stealth]
  \node[c20n=cgreen] (o) at (0,0) {open()};
  \node[c20n=cblue] (b) at (4.2,0) {your code runs};
  \node[c20n=cred] (c) at (8.4,0) {close()};
  \draw[->,thick] (o) -- (b);
  \draw[->,thick] (b) -- node[above,font=\sffamily\scriptsize]{block ends} (c);
  \draw[->,thick,cgold,rounded corners=3pt] (b.south) -- ++(0,-0.7) -| node[pos=0.25,below,font=\sffamily\scriptsize,text=ink]{or an error is raised} (c.south);
  \draw[decorate,decoration={brace,mirror,amplitude=4pt},cgray] (-0.9,-1.85) -- node[below=4pt,font=\sffamily\scriptsize,text=ink]{\code{with open(...) as f:} manages all of this} (9.4,-1.85);
\end{tikzpicture}
\end{center}
% MD-BEGIN
% ```mermaid
% flowchart LR
%     O["open()"] --> B["your code runs"]
%     B -- "block ends" --> C["close()"]
%     B -- "or an error is raised" --> C
% ```
% *`with open(...) as f:` manages all of this.*
% MD-END

\begin{pycode}
with open("notes.txt", "a", encoding="utf-8", newline="\n") as f:
    print("line four", file=f)
    f.write("line five\n")
print(f.closed)
with open("notes.txt", encoding="utf-8") as f:
    text = f.read()
print(type(text), len(text))
print(text, end="")
\end{pycode}
\begin{rawcode}
True
<class 'str'> 49
line one
line two
line three
line four
line five
\end{rawcode}

Two ways to write text: \code{f.write(s)} writes exactly \code{s}. \code{print(s, file=f)} is
\code{print} aimed at a file, so it adds a newline unless you pass \code{end=""}. It also joins
several arguments with spaces unless you change \code{sep}.

\subsection{Read a text file}

\begin{center}
\begin{tabularx}{\linewidth}{@{}l l X@{}}
\toprule
\textbf{Call} & \textbf{Returns} & \textbf{Memory use} \\
\midrule
\code{f.read()} & one \code{str}, the whole file & whole file in memory \\
\code{f.read(n)} & next \code{n} characters & small, you choose \\
\code{f.readline()} & one line, with its \code{\textbackslash n} & one line \\
\code{f.readlines()} & \code{list} of all lines & whole file in memory \\
\code{for line in f:} & one line at a time & one line; best for big files \\
\bottomrule
\end{tabularx}
\end{center}

At the end of the file, \code{read()} and \code{readline()} return an empty string \code{""}.
That is how you detect the end. A blank line in the middle is \code{"\textbackslash n"}, which is
not empty, so it is still truthy.

\begin{pycode}
with open("notes.txt", encoding="utf-8") as f:
    print(repr(f.readline()))
    print(repr(f.read(4)))
    print(repr(f.read()))
    print(repr(f.read()), repr(f.readline()))
with open("notes.txt", encoding="utf-8") as f:
    lines = f.readlines()
print(len(lines), repr(lines[0]), repr(lines[-1]))
with open("notes.txt", encoding="utf-8") as f:
    for i, line in enumerate(f, 1):
        print(i, line.rstrip("\n"))
\end{pycode}
\begin{rawcode}
'line one\n'
'line'
' two\nline three\nline four\nline five\n'
'' ''
5 'line one\n' 'line five\n'
1 line one
2 line two
3 line three
4 line four
5 line five
\end{rawcode}

\begin{typecheck}
A file object is its own iterator: \code{iter(f) is f}. Looping over it yields \code{str} lines
\emph{including} the trailing newline, so you often write \code{line.rstrip("\textbackslash n")}.
Because it is an iterator, you can use it once. A second loop over the same open file yields
nothing, since the pointer is already at the end.
\end{typecheck}

\begin{gamerule}
For anything that could be large (logs, datasets), loop over the file object. Do not call
\code{read()} or \code{readlines()}: a 5 GB file would need 5 GB of memory.
\end{gamerule}

\subsubsection*{Example: JSONL, one JSON object per line}

JSONL is the common format for eval sets and fine-tuning data: each line is a complete JSON
document. Write with \code{json.dumps(obj) + "\textbackslash n"}, read with \code{json.loads(line)}.

\begin{pycode}
import json

rows = [{"id": 1, "score": 0.5}, {"id": 2, "score": 0.75}]
with open("out.jsonl", "w", encoding="utf-8", newline="\n") as f:
    for r in rows:
        f.write(json.dumps(r) + "\n")
with open("out.jsonl", "rb") as f:
    print(f.read())
with open("out.jsonl", encoding="utf-8") as f:
    back = [json.loads(line) for line in f]
print(back == rows, type(back[0]))
with open("out.jsonl", encoding="utf-8") as f:
    print(sum(json.loads(line)["score"] for line in f if line.strip()))
\end{pycode}
\begin{rawcode}
b'{"id": 1, "score": 0.5}\n{"id": 2, "score": 0.75}\n'
True <class 'dict'>
1.25
\end{rawcode}

The last loop streams the file and keeps one line in memory at a time, so it scales to files much
bigger than RAM.

\subsection{Errors you will meet}

\begin{pycode}
try:
    open("notes.txt", "x", encoding="utf-8")
except FileExistsError as e:
    print("FileExistsError", e.errno)
try:
    open("missing.txt", encoding="utf-8")
except FileNotFoundError as e:
    print("FileNotFoundError:", e.strerror)
with open("cafe.txt", "w", encoding="utf-8") as f:
    f.write("caf\u00e9\n")
try:
    with open("cafe.txt", encoding="ascii") as f:
        f.read()
except UnicodeDecodeError as e:
    print("UnicodeDecodeError:", e.reason)
\end{pycode}
\begin{rawcode}
FileExistsError 17
FileNotFoundError: No such file or directory
UnicodeDecodeError: ordinal not in range(128)
\end{rawcode}

The error text for a missing file differs between operating systems, but the exception class
does not. Catch the class. The third error is the str/bytes border from the Text Data chapter
again: you told Python the bytes were ASCII, but the file is UTF-8.

\subsection{Newlines}

Lines end with \code{"\textbackslash n"} on Linux and macOS and \code{"\textbackslash r\textbackslash n"} on Windows. In text mode Python translates for you:
reading turns any ending into \code{"\textbackslash n"}, and writing turns \code{"\textbackslash n"} into the platform's ending
unless you pass \code{newline}.

\begin{pycode}
with open("a.txt", "w", newline="") as f:        # no translation
    f.write("a\nb\n")
with open("b.txt", "w", newline="\r\n") as f:    # force Windows endings
    f.write("a\nb\n")
print(open("a.txt", "rb").read(), open("b.txt", "rb").read())
print(open("b.txt", encoding="utf-8").read().splitlines())
\end{pycode}
\begin{rawcode}
b'a\nb\n' b'a\r\nb\r\n'
['a', 'b']
\end{rawcode}

\subsection{Binary files and seek()}

Add \code{b} to the mode and you get \code{bytes} back, with no decoding and no newline
translation. Use it for images, model weights, pickles and anything that is not text.

\begin{pycode}
data = bytes(range(256))
with open("b.bin", "wb") as f:
    print(f.write(data))
with open("b.bin", "rb") as f:
    raw = f.read()
print(type(raw), len(raw), raw[:4])
\end{pycode}
\begin{rawcode}
256
<class 'bytes'> 256 b'\x00\x01\x02\x03'
\end{rawcode}

A file object remembers a \emph{position} (a pointer). Every \code{read} and \code{write} starts at
the pointer and moves it forward. \code{tell()} reports it, \code{seek(offset, origin)} moves it.
The position counts \emph{bytes} from the start of the file.

\begin{center}
\begin{tikzpicture}[>=Stealth]
  \foreach \i in {0,...,9} \node[c20c] (c\i) at (\i*0.75,0) {\i};
  \node[font=\sffamily\scriptsize,anchor=east] at (-0.5,0) {byte offset};
  \node[font=\sffamily\scriptsize,text=cgray] at (6.75,0.6) {...};
  % pointer after read(3)
  \draw[->,thick,cgreen] (2.625,-1.0) -- (2.625,-0.4);
  \node[font=\ttfamily\scriptsize,text=cgreen!60!black,anchor=north] at (2.625,-1.0) {after read(3): tell() = 3};
  \node[font=\sffamily\scriptsize,text=cgray,anchor=west] at (7.4,0.0) {file length N};
\end{tikzpicture}
\end{center}
% MD-BEGIN
% ```text
% offset :   0   1   2   3   4   5   6   7   8   9  ...
% bytes  : [00][01][02][03][04][05][06][07][08][09] ...
%                         ^
%                         pointer after read(3)  -> tell() == 3
% ```
% MD-END

\begin{center}
\begin{tikzpicture}[>=Stealth]
  \foreach \i in {0,...,9} \node[c20c] (c\i) at (\i*0.75,0) {};
  \node[font=\sffamily\scriptsize,anchor=east] at (-0.5,0) {file};
  \node[font=\ttfamily\scriptsize,text=cblue] at (0,0.75) {start};
  \node[font=\ttfamily\scriptsize,text=cred] at (6.75,0.75) {end};
  \draw[->,thick,cblue] (0,-0.5) -- node[below,font=\ttfamily\scriptsize]{seek(3, 0)} (2.25,-0.5);
  \draw[->,thick,cgreen] (2.25,-1.2) -- node[below,font=\ttfamily\scriptsize]{seek(2, 1)} (3.75,-1.2);
  \draw[->,thick,cred] (6.75,-1.9) -- node[below,font=\ttfamily\scriptsize]{seek(-3, 2)} (4.5,-1.9);
  \node[font=\sffamily\scriptsize,anchor=west,text=cblue] at (7.5,-0.5) {0 = from the start};
  \node[font=\sffamily\scriptsize,anchor=west,text=cgreen!60!black] at (7.5,-1.2) {1 = from the current position};
  \node[font=\sffamily\scriptsize,anchor=west,text=cred] at (7.5,-1.9) {2 = from the end (offset is negative)};
\end{tikzpicture}
\end{center}
% MD-BEGIN
% ```text
%         start                                    end
% offset:   0  1  2  3  4  5  6  7  8  9 (N=10)
%
% seek(3, 0)   -> position 3          0 = from the start (os.SEEK_SET)
% seek(2, 1)   -> current + 2         1 = from the current position (os.SEEK_CUR)
% seek(-3, 2)  -> end - 3 = 7         2 = from the end, offset negative (os.SEEK_END)
% ```
% MD-END

\begin{pycode}
import os

with open("b.bin", "rb") as f:
    print(f.tell(), f.seek(255), f.read(), f.tell())
    print(f.seek(-1, os.SEEK_END), f.read(1))
    print(f.seek(10), f.seek(5, os.SEEK_CUR), f.read(3))
print(os.SEEK_SET, os.SEEK_CUR, os.SEEK_END)
\end{pycode}
\begin{rawcode}
0 255 b'\xff' 256
255 b'\xff'
10 15 b'\x0f\x10\x11'
0 1 2
\end{rawcode}

\code{seek} returns the new position. Reading the last byte of a huge file this way is instant,
because nothing before it is read.

\begin{trap}
In text mode, \code{seek} and \code{tell} work with opaque numbers. Only \code{seek(0)} and values
that \code{tell()} gave you are safe. With UTF-8, one character can be several bytes, so ``go to
character 100'' is not ``go to byte 100''. For arbitrary jumps, open the file in binary mode.
\end{trap}

\begin{boss}[: the file pointer]
A file object is not the file. It is a cursor plus a buffer. Reading moves the cursor, so the
\emph{same} file object returns different data on each call, and returns nothing once it is at
the end. Writing in mode \code{"w"} erases first, then moves forward. Most confusing file bugs
come from forgetting where the cursor is: a second \code{read()} that returns \code{""}, a loop
over a file that runs only once, or a write that overwrote the start because of \code{"r+"} and
no \code{seek}. When in doubt, print \code{f.tell()}.
\end{boss}

\subsection{Memory mapping}

The standard \code{mmap} module can map a file into memory so that it looks like a
\code{bytearray}: you slice it and search it, and the operating system loads pages on demand.
It helps for very large binary files you need random access to. For everyday work, plain
\code{open} is enough.

% ------------------------------------------------------------
\section{File Operations}

Python's names come from Unix. The \code{os}, \code{os.path} and \code{shutil} modules hold the
functions. The newer \code{pathlib} module (later in this chapter) does the same jobs with
objects.

\begin{center}
\begin{tabularx}{\linewidth}{@{}l l X@{}}
\toprule
\textbf{Task} & \textbf{Function} & \textbf{Returns / note} \\
\midrule
exists? & \code{os.path.exists(p)} & \code{bool}; works for files and directories \\
is a file? & \code{os.path.isfile(p)} & \code{bool} \\
is a directory? & \code{os.path.isdir(p)} & \code{bool} \\
is it absolute? & \code{os.path.isabs(p)} & \code{bool}; does not check the disk \\
size & \code{os.path.getsize(p)} & \code{int} bytes \\
copy & \code{shutil.copy(src, dst)} & returns the destination path \\
move or rename & \code{os.rename(a, b)}, \code{shutil.move} & \code{move} also works across disks \\
delete a file & \code{os.remove(p)} & permanent; no recycle bin \\
hard / soft link & \code{os.link}, \code{os.symlink} & Unix concepts; \code{islink()} tests \\
permissions & \code{os.chmod(p, mode)} & Unix; \code{0o400} = owner can read only \\
owner & \code{os.chown(p, uid, gid)} & Unix only \\
\bottomrule
\end{tabularx}
\end{center}

\begin{pycode}
import shutil

with open("oops.txt", "w", encoding="utf-8", newline="\n") as f:
    f.write("hello\n")
print(os.path.exists("oops.txt"), os.path.exists("waffles"))
print(os.path.exists("."))
print(os.path.isfile("oops.txt"), os.path.isdir("oops.txt"), os.path.isdir("."))
print(os.path.isabs("oops.txt"), os.path.isabs(os.path.abspath("oops.txt")))
print(shutil.copy("oops.txt", "copy.txt"))
os.rename("copy.txt", "renamed.txt")
print(sorted(x for x in os.listdir(".") if x.endswith(".txt")))
os.remove("renamed.txt")
print(os.path.exists("renamed.txt"), os.path.getsize("oops.txt"))
\end{pycode}
\begin{rawcode}
True False
True
True False True
False True
copy.txt
['a.txt', 'b.txt', 'cafe.txt', 'notes.txt', 'oops.txt', 'renamed.txt']
False 6
\end{rawcode}

\emph{Links.} A \emph{hard link} is a second name for the same bytes on disk. A \emph{symbolic link}
is a small file that points to another path. \code{os.symlink} needs special rights on Windows,
so it is not run here. \code{chmod} and \code{chown} are for Unix-like systems. A typical call is
\code{os.chmod("secrets.txt", 0o600)} so that only you can read your key file.

\begin{trap}
\code{os.remove} and \code{shutil.rmtree} delete for good. \code{os.rename} silently
overwrites an existing destination on Unix. Check with \code{os.path.exists} first, or write to a
temporary name and then \code{Path.replace()} into place, which swaps the file in one step.
\end{trap}

% ------------------------------------------------------------
\section{Directory Operations}

\begin{center}
\begin{tabularx}{\linewidth}{@{}l l X@{}}
\toprule
\textbf{Task} & \textbf{Function} & \textbf{Note} \\
\midrule
create & \code{os.mkdir(d)} & one level; \code{FileExistsError} if present \\
create nested & \code{os.makedirs(d, exist\_ok=True)} & all missing levels, no error if present \\
delete empty & \code{os.rmdir(d)} & fails if not empty \\
delete tree & \code{shutil.rmtree(d)} & deletes everything inside: careful \\
list names & \code{os.listdir(d)} & \code{list} of \code{str}, unsorted \\
current dir & \code{os.getcwd()}, \code{os.chdir(d)} & changes for the whole program \\
match names & \code{glob.glob(pattern)} & shell-style patterns, not regex \\
\bottomrule
\end{tabularx}
\end{center}

\begin{pycode}
import glob

os.makedirs("data/runs/a1", exist_ok=True)
print(os.listdir("data"), os.path.isdir("data/runs/a1"))
try:
    os.rmdir("data")
except OSError as e:
    print("OSError:", type(e).__name__)
shutil.rmtree("data")
print(os.path.exists("data"))
os.mkdir("runs")
for name in ("a.jsonl", "b.jsonl", "notes.md"):
    with open("runs/" + name, "w", encoding="utf-8") as f:
        f.write("x\n")
os.chdir("runs")
print(sorted(os.listdir(".")), os.path.basename(os.getcwd()))
os.chdir("..")
\end{pycode}
\begin{rawcode}
['runs'] True
OSError: OSError
False
['a.jsonl', 'b.jsonl', 'notes.md'] runs
\end{rawcode}

The glob patterns are \emph{not} regular expressions.

\begin{center}
\begin{tabularx}{\linewidth}{@{}l X l@{}}
\toprule
\textbf{Pattern} & \textbf{Matches} & \textbf{Regex equivalent} \\
\midrule
\code{*} & any number of any characters & \code{.*} \\
\code{?} & exactly one character & \code{.} \\
\code{[abc]} & one of \code{a}, \code{b}, \code{c} & same \\
\code{[!abc]} & one character that is not \code{a}, \code{b}, \code{c} & \code{[\^{}abc]} \\
\bottomrule
\end{tabularx}
\end{center}

\begin{pycode}
print(sorted(glob.glob("*.txt")))
print(glob.glob("ru??"), glob.glob("[rz]*"))
print(sorted(os.path.basename(p) for p in glob.glob("runs/*.jsonl")))
\end{pycode}
\begin{rawcode}
['a.txt', 'b.txt', 'cafe.txt', 'notes.txt', 'oops.txt']
['runs'] ['runs']
['a.jsonl', 'b.jsonl']
\end{rawcode}

\begin{trap}
\code{os.listdir} gives names only, not paths, and in no fixed order. Wrap it in \code{sorted()}
if order matters, and join the folder name back on before opening a file. Also \code{os.chdir}
changes the whole program's idea of ``here''; avoid it in library code and use full paths.
\end{trap}

% ------------------------------------------------------------
\section{Pathnames}

A \emph{pathname} says where a file is: \emph{absolute} from the root, or \emph{relative} to the
current directory. Unix, macOS and URLs separate parts with \code{/}. Windows uses
\code{\textbackslash}. Python accepts \code{/} on Windows too, so you can often use slashes
everywhere. If you write backslashes yourself, double them or use a raw string
(\code{r"eek\textbackslash urk"}); \code{"eek\textbackslash urk"} contains an escape sequence, not a path.

\subsection{os.path}

\begin{pycode}
p = os.path.join("eek", "urk", "snort.txt")        # uses the OS separator
print(p.replace(os.sep, "/"))
print(os.path.splitext("a/b/model.tar.gz"))
print(os.path.basename("a/b/c.txt"), os.path.dirname("a/b/c.txt"))
print(os.path.isabs(os.path.abspath("x")))
\end{pycode}
\begin{rawcode}
eek/urk/snort.txt
('a/b/model.tar', '.gz')
c.txt a/b
True
\end{rawcode}

\code{os.path.abspath(p)} turns a relative path into an absolute one. \code{os.path.realpath(p)}
also follows symbolic links to the real file. Never build paths with \code{"a" + "/" + "b"}: let
\code{join} or \code{pathlib} insert the right separator.

\subsection{pathlib}

\code{pathlib.Path} treats a path as an \emph{object}, not a string. The \code{/} operator joins
parts, and the object has handy attributes and methods.

\begin{center}
\begin{tikzpicture}[>=Stealth]
  \foreach \x/\t/\c in {0/{runs}/cblue, 2.0/{/}/cgray, 2.8/{2025}/cblue, 4.8/{/}/cgray, 5.6/{eval.jsonl}/cgreen}
    \node[font=\ttfamily\Large,text=\c] at (\x,0) {\t};
  \draw[decorate,decoration={brace,amplitude=4pt},cblue,thick] (-0.5,0.45) -- node[above=3pt,font=\ttfamily\scriptsize,text=ink]{.parent = runs/2025} (4.0,0.45);
  \draw[decorate,decoration={brace,mirror,amplitude=4pt},cgreen,thick] (4.7,-0.45) -- node[below=3pt,font=\ttfamily\scriptsize,text=ink]{.name = eval.jsonl} (7.1,-0.45);
  \draw[decorate,decoration={brace,mirror,amplitude=3pt},cgold,thick] (4.7,-1.3) -- node[below=3pt,font=\ttfamily\scriptsize,text=ink]{.stem = eval} (6.2,-1.3);
  \draw[decorate,decoration={brace,mirror,amplitude=3pt},cpurple,thick] (6.25,-1.3) -- node[below=3pt,font=\ttfamily\scriptsize,text=ink]{.suffix} (7.1,-1.3);
\end{tikzpicture}
\end{center}
% MD-BEGIN
% ```text
% Path("runs/2025/eval.jsonl")
%       |-------- .parent --------|  |------ .name ------|
%       runs / 2025                  /  eval  .jsonl
%                                       .stem  .suffix
% .parts = ('runs', '2025', 'eval.jsonl')
% ```
% MD-END

\begin{pycode}
from pathlib import Path, PurePosixPath, PureWindowsPath

p = Path("eek") / "urk" / "snort.txt"
print(p.as_posix(), p.name, p.stem, p.suffix)
print(p.parent.as_posix(), p.parts)
print(p.with_suffix(".jsonl").as_posix(), p.with_name("x.md").as_posix())
print(PureWindowsPath(PurePosixPath("eek/urk/snort.txt")))
print(type(p).__name__ in ("WindowsPath", "PosixPath"))
\end{pycode}
\begin{rawcode}
eek/urk/snort.txt snort.txt snort .txt
eek/urk ('eek', 'urk', 'snort.txt')
eek/urk/snort.jsonl eek/urk/x.md
eek\urk\snort.txt
True
\end{rawcode}

\code{print(p)} shows the path in your platform's style, and \code{repr(p)} shows
\code{WindowsPath(...)} or \code{PosixPath(...)} depending on the system. The \code{Pure...}
classes only manipulate path text and never touch the disk, so you can build a Windows path on a
Mac.

\code{Path} objects also do the file work:

\begin{pycode}
runs = Path("runs")
body = '{"x": 2}\n{"x": 3}\n'
(runs / "c.jsonl").write_text(body, encoding="utf-8", newline="\n")
print(sorted(x.name for x in runs.iterdir()))
print(sorted(x.name for x in runs.glob("*.jsonl")))
print(sorted(x.as_posix() for x in Path(".").rglob("*.md")))
print((runs / "c.jsonl").read_text(encoding="utf-8").splitlines())
print((runs / "c.jsonl").exists(), (runs / "c.jsonl").is_file(), runs.is_dir())
print((runs / "c.jsonl").stat().st_size, runs.resolve().is_absolute())
tmp = Path("final.json.tmp")
tmp.write_text("{}", encoding="utf-8")
tmp.replace("final.json")
print(sorted(x for x in os.listdir(".") if x.startswith("final")))
\end{pycode}
\begin{rawcode}
['a.jsonl', 'b.jsonl', 'c.jsonl', 'notes.md']
['a.jsonl', 'b.jsonl', 'c.jsonl']
['runs/notes.md']
['{"x": 2}', '{"x": 3}']
True True True
18 True
['final.json']
\end{rawcode}

\begin{center}
\begin{tabularx}{\linewidth}{@{}l l l@{}}
\toprule
\textbf{Job} & \textbf{os / os.path style} & \textbf{pathlib style} \\
\midrule
join & \code{os.path.join(a, b)} & \code{Path(a) / b} \\
file name & \code{os.path.basename(p)} & \code{p.name} \\
extension & \code{os.path.splitext(p)[1]} & \code{p.suffix} \\
exists & \code{os.path.exists(p)} & \code{p.exists()} \\
read all text & \code{open(p).read()} & \code{p.read\_text(encoding=...)} \\
list folder & \code{os.listdir(d)} & \code{d.iterdir()} \\
find by pattern & \code{glob.glob("d/*.py")} & \code{d.glob("*.py")} \\
\bottomrule
\end{tabularx}
\end{center}

\begin{gamerule}
Prefer \code{pathlib} in new code. Paths become objects with a clear type, the \code{/} join is
portable, and one-line helpers like \code{read\_text} and \code{write\_text} close the file for
you. You can pass a \code{Path} anywhere a filename string is accepted, including \code{open()}.
\end{gamerule}

% ------------------------------------------------------------
\section{BytesIO and StringIO}

Many functions want a file object, but your data is already in memory. \code{io.StringIO} wraps a
\code{str} and \code{io.BytesIO} wraps \code{bytes}. They have the same methods as real files
(\code{read}, \code{write}, \code{seek}, iteration) but no disk is touched.

\begin{center}
\begin{tikzpicture}[>=Stealth]
  \node[c20n=cblue,align=center] (fn) at (0,0) {function that wants\\[-1pt]a file object};
  \node[c20n=cgray] (real) at (6.5,0.9) {open("x.txt")};
  \node[c20n=cgreen] (sio) at (6.5,0) {StringIO("text")};
  \node[c20n=cpurple] (bio) at (6.5,-0.9) {BytesIO(b"bytes")};
  \draw[<-,thick,cgray] (fn.east) -- (real.west);
  \draw[<-,thick,cgray] (fn.east) -- (sio.west);
  \draw[<-,thick,cgray] (fn.east) -- (bio.west);
  \node[font=\sffamily\scriptsize,text=cgray,anchor=west] at (9.2,0.9) {disk};
  \node[font=\sffamily\scriptsize,text=cgray,anchor=west] at (9.2,0) {memory, text};
  \node[font=\sffamily\scriptsize,text=cgray,anchor=west] at (9.2,-0.9) {memory, bytes};
\end{tikzpicture}
\end{center}
% MD-BEGIN
% ```mermaid
% flowchart LR
%     R["open('x.txt') (disk)"] --> F["function that wants a file object"]
%     S["StringIO('text') (memory, text)"] --> F
%     B["BytesIO(b'bytes') (memory, bytes)"] --> F
% ```
% MD-END

\begin{pycode}
import io, csv

s = io.StringIO()
s.write("a,b\n")
print("x", file=s)
print(repr(s.getvalue()), type(s).__name__)
s.seek(0)
print(s.readline(), end="")
print(list(io.StringIO("p\nq\n")))
bio = io.BytesIO(b"\x89PNG\r\n")
print(bio.read(4), bio.tell(), bio.getvalue())
bio.seek(0)
bio.write(b"AB")
print(bio.getvalue())
print(list(csv.DictReader(io.StringIO("name,score\nx,1\ny,2\n"))))
\end{pycode}
\begin{rawcode}
'a,b\nx\n' StringIO
a,b
['p\n', 'q\n']
b'\x89PNG' 4 b'\x89PNG\r\n'
b'ABNG\r\n'
[{'name': 'x', 'score': '1'}, {'name': 'y', 'score': '2'}]
\end{rawcode}

\begin{typecheck}
\code{StringIO} holds \code{str} and \code{BytesIO} holds \code{bytes}. Do not mix them: writing a
\code{str} into a \code{BytesIO} raises \code{TypeError}. \code{getvalue()} returns the whole
content regardless of the pointer, while \code{read()} starts from the pointer. After you write
into a \code{StringIO}, call \code{seek(0)} before reading it back.
\end{typecheck}

Typical uses: unit tests for code that reads a file, capturing what a function prints into a
string, building a CSV or an image in memory before sending it over the network, and converting an
image format with the Pillow library (\code{Image.open(BytesIO(data))}, which needs
\code{pip install pillow}; not run here).

% ------------------------------------------------------------
\section{File Formats: Determination}

You may meet a file and not know what it is. The file extension is only a hint: anyone can
rename \code{.png} to \code{.txt}. Better checks look at the first bytes, the \emph{magic number}
(chapter on binary data: a PNG starts with the bytes \code{\textbackslash x89PNG}).

\begin{center}
\begin{tabularx}{\linewidth}{@{}l X@{}}
\toprule
\textbf{Tool} & \textbf{What it does} \\
\midrule
\code{file} (Unix command) & identifies files by magic bytes and other rules \\
\code{python-magic} & Python wrapper around the \code{libmagic} library used by \code{file} \\
\code{magika} & Google's learned file-type detector; \code{pip install magika} \\
your own check & compare \code{f.read(8)} with a known header such as \code{b"\textbackslash x89PNG\textbackslash r\textbackslash n\textbackslash x1a\textbackslash n"} \\
\bottomrule
\end{tabularx}
\end{center}

\begin{pycode}
sig = b"\x89PNG\r\n\x1a\n"
with open("fake.png", "wb") as f:
    f.write(sig + b"rest of image")
with open("fake.png", "rb") as f:
    print(f.read(8) == sig)
\end{pycode}
\begin{rawcode}
True
\end{rawcode}

% ------------------------------------------------------------
\section{Review and what comes next}

You can now open files in the right mode, read them line by line without loading them whole, use
\code{with} so that they close, move the file pointer with \code{seek}, and manage files and
folders with \code{os}, \code{shutil} and \code{pathlib}. The next chapter is about concurrency:
doing several things at once.

\begin{practice}
\begin{enumerate}[nosep,leftmargin=*]
  \item Write the string \code{"This is a test of the emergency text system"} to \code{test.txt}
        with a \code{with} block and UTF-8. Read it back into \code{test2}. Are they equal?
  \item List all \code{.txt} files in the current folder, sorted, using \code{pathlib} and then
        using \code{glob}.
  \item Why does a second \code{for line in f:} over the same open file print nothing? Show two
        ways to fix it.
  \item You have a 4 GB JSONL file of eval rows. Write the loop that counts rows with
        \code{"score" > 0.5} without loading the file into memory.
  \item Write 100 bytes, then use \code{seek} to read the last 10 bytes in two different ways.
  \item What goes wrong if you open an existing file with mode \code{"w"} just to look at it?
        Which mode would you use to append a line to a log?
\end{enumerate}
\end{practice}

\begin{achievement}{File Handler}
\begin{itemize}[nosep,leftmargin=*]
  \item Choose the correct open mode and always close files with \code{with}.
  \item Stream large text files line by line and read or write JSONL.
  \item Move around a binary file with \code{seek} and \code{tell}.
  \item Build, split and search paths with \code{pathlib} and \code{glob}.
  \item Use \code{StringIO} and \code{BytesIO} to treat memory like a file.
\end{itemize}
\end{achievement}
